\VerdierPhysicalPageAnchor{0037}%
\VerdierSemanticAnchor{ch1-1-2-1-coherence}
et pour tout triplet $(g,h,k)$ d'éléments de $G$, on a la relation :

\VerdierSemanticAnchor{ch1-1-2-1-eq3}
\begin{equation}
  \bigl(T(g)\star c(k,h)\bigr)c(kh,g)
  =\bigl(c(h,g)\star T(k)\bigr)c(k,hg).
  \tag{1.2.1.3}
\end{equation}

\VerdierSemanticAnchor{ch1-1-2-2}
\noindent\textbf{Définition 1.2.2.}\quad
Soit $G$ un groupe. On appelle $G^\circ$-\emph{catégorie}
($G^\circ$ groupe opposé à $G$) une catégorie $\mathcal C_e$ munie d'une
famille de foncteurs $T(g)$ ($g\in G$) et d'isomorphismes de foncteurs
$c(h,g)$ ($g,h\in G$) vérifiant les propriétés (1.2.1.1), (1.2.1.2) et
(1.2.1.3).

\VerdierSemanticAnchor{ch1-1-2-2-transition}
Une $G^\circ$-catégorie est donc une catégorie $\mathcal C_e$ munie d'une
opération à isomorphisme près d'un groupe $G^\circ$. On utilise la notation
$(\mathcal C_e,G^\circ,T,c)$ pour désigner la $G^\circ$-catégorie ; les
isomorphismes $c(h,g)$ sont appelés les
\emph{isomorphismes de transition}.

\VerdierSemanticAnchor{ch1-1-2-2-strict}
Lorsque les isomorphismes de transition sont des isomorphismes identiques,
ce qui implique en particulier qu'on a pour tout couple $g,h\in G$
l'égalité $T(hg)=T(g)T(h)$, on dit que la $G^\circ$-catégorie est
\emph{stricte}. Un tel objet est donc défini par une représentation de
$G^\circ$ dans le groupe des automorphismes de la catégorie $\mathcal C_e$.

\VerdierSemanticAnchor{ch1-1-2-3}
\noindent\textbf{1.2.3.}\quad
Soient maintenant $\mathcal C$ et $\mathcal C'$ deux catégories graduées de
type $G$ et $F\colon\mathcal C\to\mathcal C'$ un foncteur gradué.
Choisissons sur $\mathcal C$ et $\mathcal C'$ des clivages normalisés. Le
choix de ces clivages définit des isomorphismes de foncteurs :

\VerdierSemanticAnchor{ch1-1-2-3-eq1}
\begin{equation}
  m(g)\colon F_eT(g)\longrightarrow T'(g)F_e,\qquad
  g\in G,\qquad m(e)=\text{identité de }F_e
  \tag{1.2.3.1}
\end{equation}
tels que pour tout couple $g,h\in G$, le diagramme ci-après soit
commutatif :

\VerdierSemanticAnchor{ch1-1-2-3-eq2}
\begin{equation}
\begin{tikzcd}[column sep=huge,row sep=large]
F_eT(hg)
  \arrow[r,"m(hg)"]
  \arrow[d,"{F_e\star c(h,g)}",swap]
&
T^{\prime}(hg)F_e
  \arrow[dd,"{c^{\prime}(h,g)\star F_e}"]
\\
F_eT(g)T(h)
  \arrow[d,"m(g)\star T(h)",swap]
&
{}
\\
T^{\prime}(g)F_eT(h)
  \arrow[r,"T^{\prime}(g)\star m(h)"]
&
T^{\prime}(g)T^{\prime}(h)F_e\mathpunct{.}
\end{tikzcd}
\tag{1.2.3.2}
\end{equation}

\VerdierSemanticAnchor{ch1-1-2-4}
\noindent\textbf{Définition 1.2.4.}\quad
Soient $(\mathcal C_e,G^\circ,T,c)$ et
$(\mathcal C'_e,G^\circ,T',c')$ deux $G^\circ$-catégories. On appelle
$G^\circ$-\emph{foncteur} la donnée d'un foncteur :

\VerdierSemanticAnchor{ch1-1-2-4-fe}
\[
  F_e\colon\mathcal C_e\longrightarrow\mathcal C'_e
\]
